\documentclass[10pt,a4paper]{amsart}
\usepackage[margin=19mm]{geometry}
\usepackage{amsmath,amssymb,mathtools}
\usepackage[scr=rsfs,cal=euler]{mathalfa}
\usepackage{xfrac}
\usepackage[unicode,hidelinks,bookmarks=false]{hyperref}
\newcommand{\ie}{i.e.\ }
\newcommand{\cf}{cf.\ }
\newcommand{\resp}{respectively\ }
\newcommand{\Subsection}{\subsection}
\providecommand{\nobreakdash}{\nobreak\hbox{-}}
\setlength{\emergencystretch}{2em}
\multlinegap0pt
\usepackage{bm}
\usepackage{bbm}
\usepackage{dsfont}
\usepackage{xspace}
\usepackage{cancel}
\usepackage{mathtools}
\usepackage{amsthm}
\makeatletter
\@ifundefined{Theorem}{\newtheorem{Theorem}{Theorem}}{}
\@ifundefined{Lemma}{\newtheorem{Lemma}[Theorem]{\sc Lemma}}{}
\makeatother
\title[An improved Copson inequality]{An improved Copson inequality}
\author{Bikram Das, Atanu Manna}
\date{Source: arXiv:2508.00388v1 (CC BY 4.0)}
\begin{document}
\maketitle

\noindent\emph{Source and licence.} An improved Copson inequality. Version arXiv:2508.00388v1 (CC BY 4.0), distributed under the Creative Commons Attribution 4.0 International licence, as recorded in the arXiv licence metadata for this exact version and shown on the abstract page https://arxiv.org/abs/2508.00388v1. Full rights notice: licenses/arxiv-2508.00388-CC-BY-4.0.txt.\par\smallskip

\noindent\textit{Editorial scope.} Counted unit: the Lemma whose source label is \texttt{Positive} in the article (source0.tex lines 345--347, source label \texttt{Positive}), delivered together with the article's own complete original proof of it (source0.tex lines 348--451, one \texttt{proof} environment of 104 lines). No mathematics is written, repaired or re-derived: statement and proof are the article's own text line for line. Required context retained uncounted: none. Every definition, display and label of those lines is retained unchanged. Omitted from this package and not redistributed: 1--344, 452--831. Nothing inside a retained excerpt is omitted or altered: every retained body is delivered exactly as the article's lines are written, so no omission receipt of any kind accompanies this package. Citations: the retained excerpts carry no citation command at all, so no bibliography entry of the article is transcribed and no reference is redistributed. Reference adaptation: none was needed and none was made; every reference inside the delivered unit resolves inside it, so no unresolved reference or citation is produced. Printed slips kept literal: no printed word, symbol, hypothesis or number of the counted statement or of its proof was altered. Layout and class: the article's own class is \texttt{amsart} with packages including amssymb, amsmath, amsfonts, latexsym, bm, xy, amscd; the wrapper is the lane's pinned \texttt{amsart} editorial wrapper at 10pt on A4, which restates the theorem-like environments the retained text needs and adds the article's own macro definitions that the retained text needs (none), with extra wrapper packages (bm, bbm, dsfont, xspace, cancel, mathtools). The delivered numbering is the unit's own. Assets: no retained span contains a figure environment, a graphics-inclusion command, a PDF-page inclusion, a table or a TikZ picture; no image file is copied into this package, the package declares no asset dependency and no image file is redistributed with it. Licence: version arXiv:2508.00388v1 of this article is distributed under the Creative Commons Attribution 4.0 International licence, as recorded in the arXiv licence metadata for this exact version and shown on the abstract page https://arxiv.org/abs/2508.00388v1. A full rights notice is delivered at licenses/arxiv-2508.00388-CC-BY-4.0.txt. Characters: every retained excerpt is pure ASCII, so the delivered engine, XeLaTeX with its default Latin Modern text font, reports no missing character.
\begin{Lemma}{\label{Positive}}
Let $\alpha\in[\frac{1}{3},1)$ and $\{q_{n}\}$ be decreasing sequence of positive real numbers. Then $C_{k}(q,\alpha)$ is strictly positive for $k\geq3$, $k\in\mathbb{N}$.
\end{Lemma}

\begin{proof}
Denote $\gamma_{1}=\frac{1-\alpha}{2}$ and $\gamma_{2}=\frac{1+\alpha}{2}$ for $\alpha\in[\frac{1}{3},1)$. Then one can re-write $C_{k}(q,\alpha)$ as below:
\begin{align*}
C_{k}(q,\alpha)&=(-1)^{k-1}\frac{\alpha}{k!}\frac{q_{n}q_{n+1}^{k-1}}{Q_{n}^{k}}\displaystyle\prod_{i=1}^{k-1}(i-\alpha)
+\frac{\gamma_{1}}{k!}\frac{q_{n}^{k}}{Q_{n}^{k}}\displaystyle\prod_{i=1}^{k-1}(i-\gamma_{1})+(-1)^{k}\frac{\gamma_{2}}{k!}
\frac{q_{n}q_{n+1}^{k-1}}{Q_{n}^{k}}\displaystyle\prod_{i=1}^{k-1}(i-\gamma_{2}).
\end{align*}
We now divide $k\geq 3$, $k\in\mathbb{N}$ in two cases, one is odd natural number $k$, and another is even natural number $k$.\\
\emph{\textbf{Case I: $k$ is odd.}}\\
Note that when $k$ is odd then
\begin{align*}
C_{k}(q,\alpha)&=\frac{\alpha}{k!}\frac{q_{n}q_{n+1}^{k-1}}{Q_{n}^{k}}\displaystyle\prod_{i=1}^{k-1}(i-\alpha)
+\frac{\gamma_{1}}{k!}\frac{q_{n}^{k}}{Q_{n}^{k}}\displaystyle\prod_{i=1}^{k-1}(i-\gamma_{1})-\frac{\gamma_{2}}{k!}
\frac{q_{n}q_{n+1}^{k-1}}{Q_{n}^{k}}\displaystyle\prod_{i=1}^{k-1}(i-\gamma_{2})\\
&=I_{1}+I_{2}-I_{3},
\end{align*}where
$I_{1}=\frac{\alpha}{k!}\frac{q_{n}q_{n+1}^{k-1}}{Q_{n}^{k}}\displaystyle\prod_{i=1}^{k-1}(i-\alpha)$,
$I_{2}=\frac{\gamma_{1}}{k!}\frac{q_{n}^{k}}{Q_{n}^{k}}\displaystyle\prod_{i=1}^{k-1}(i-\gamma_{1})$, and
$I_{3}=\frac{\gamma_{2}}{k!}\frac{q_{n}q_{n+1}^{k-1}}{Q_{n}^{k}}\displaystyle\prod_{i=1}^{k-1}(i-\gamma_{2}).$\\
Now for odd $k\geq3$, and $\alpha\in[\frac{1}{3},1)$, it is immediate that $I_{1}>0$. But the following difference gives
\begin{align*}
I_{2}-I_{3}&=\frac{\gamma_{2}}{k!}\frac{q_{n}^{k}}{Q_{n}^{k}}\displaystyle\prod_{i=1}^{k-1}(i-\gamma_{2})\Big[\frac{\gamma_{1}}{\gamma_{2}}\displaystyle\prod_{i=1}^{k-1}
\big(\frac{i-\gamma_{1}}{i-\gamma_{2}}\big)-\frac{q_{n+1}^{k-1}}{q_{n}^{k-1}}\Big]\\
&=\frac{\gamma_{2}}{k!}\frac{q_{n}^{k}}{Q_{n}^{k}}\displaystyle\prod_{i=1}^{k-1}(i-\gamma_{2})L(q,\alpha),
\end{align*}where we denote $L(q,\alpha)=\frac{\gamma_{1}}{\gamma_{2}}\displaystyle\prod_{i=1}^{k-1}
\big(\frac{i-\gamma_{1}}{i-\gamma_{2}}\big)-\frac{q_{n+1}^{k-1}}{q_{n}^{k-1}}=\displaystyle\prod_{i=2}^{k-1}
\big(\frac{i-\gamma_{1}}{i-\gamma_{2}}\big)-\frac{q_{n+1}^{k-1}}{q_{n}^{k-1}}$. Since $\gamma_2>\gamma_1$, and $\{q_n\}$ is decreasing, so we have
\begin{align*}
&L(q,\alpha)>1-\frac{q_{n+1}^{k-1}}{q_{n}^{k-1}}=\frac{q_{n}^{k-1}-q_{n+1}^{k-1}}{q_{n}^{k-1}}\geq 0,
\end{align*}
which shows that $L(q,\alpha)>0$. Hence $I_{2}-I_{3}>0$. Therefore $C_{k}(q,\alpha)>0$ for all odd $k\geq3$, $k\in\mathbb{N}$.\\
\emph{\textbf{Case II: $k$ is even.}}\\
In this case, using the similar expressions for $I_1$, $I_2$, and $I_3$ in the earlier case, $C_{k}(q,\alpha)$ can be written as follows:
\begin{align*}
C_{k}(q,\alpha)&=-\frac{\alpha}{k!}\frac{q_{n}q_{n+1}^{k-1}}{Q_{n}^{k}}\displaystyle\prod_{i=1}^{k-1}(i-\alpha)+
\frac{\gamma_{1}}{k!}\frac{q_{n}^{k}}{Q_{n}^{k}}\displaystyle\prod_{i=1}^{k-1}(i-\gamma_{1})+
\frac{\gamma_{2}}{k!}\frac{q_{n}q_{n+1}^{k-1}}{Q_{n}^{k}}\displaystyle\prod_{i=1}^{k-1}(i-\gamma_{2})\\
&=-I_{1}+I_{2}+I_{3}\\
&=\frac{\gamma_{1}}{k!}\Big(\displaystyle\prod_{i=1}^{k-1}(i-\alpha)\Big)\Big[\frac{q_{n}^{k}}{Q_{n}^{k}}\displaystyle\prod_{i=1}^{k-1}\big(\frac{i-\gamma_{1}}{i-\alpha}\big)-
\frac{\alpha}{\gamma_{1}}\frac{q_{n}q_{n+1}^{k-1}}{Q_{n}^{k}}\Big]+I_{3}\\
&=\frac{\gamma_{1}}{k!}\Big(\displaystyle\prod_{i=1}^{k-1}(i-\alpha)\Big)M(q,\alpha)+I_{3},
\end{align*}where we denote $M(q,\alpha)=\frac{q_{n}^{k}}{Q_{n}^{k}}\displaystyle\prod_{i=1}^{k-1}\big(\frac{i-\gamma_{1}}{i-\alpha}\big)-\frac{\alpha}{\gamma_{1}}\frac{q_{n}q_{n+1}^{k-1}}{Q_{n}^{k}}$. Since $\frac{i-\gamma_{1}}{i-\alpha}\geq1$ for $\alpha\in[1/3,1)$ and $\{q_{n}\}$ is decreasing, so we have
\begin{align*}
M(q,\alpha)& \geq\frac{q_{n}^{k}M_{1}(\alpha)}{\gamma_{1}Q_{n}^{k}\displaystyle\prod_{i=1}^{7}(i-\alpha)}.
\end{align*}
where we choose $M_{1}(\alpha)$ as below
\begin{align*}
M_{1}(\alpha)&=\Big(\gamma_{1}\displaystyle\prod_{i=1}^{7}(i-\gamma_{1})-\alpha\displaystyle\prod_{i=1}^{7}(i-\alpha)\Big)\\
&=\Big(\frac{1-\alpha}{2}\displaystyle\prod_{i=1}^{7}(i-\frac{1-\alpha}{2})-\alpha\displaystyle\prod_{i=1}^{7}(i-\alpha)\Big)~~~(\because~~ \gamma_{1}=\frac{1-\alpha}{2}).
\end{align*}Simplify the terms of the above expression for $M_{1}(\alpha)$, we get
\begin{align}{\label{M-H-eq}}
M_{1}(\alpha)&=\frac{(1-\alpha)(3\alpha-1)(13-\alpha)}{2^{8}}H(\alpha)
\end{align}
where $H(\alpha)=85\alpha^{5}-1187\alpha^{4}+8654\alpha^{3}-24898\alpha^{2}+46941\alpha-10395$. Hence we have
\begin{align*}
H'(\alpha)&=425\alpha^{4}-4748\alpha^{3}+25962\alpha^{2}-49796\alpha+46941,\\
H''(\alpha)&=1700\alpha^{3}-14244\alpha^{2}+51924\alpha-49796,\\
H'''(\alpha)&=12(425\alpha^{2}-2374\alpha+4327).
\end{align*}
Clearly, $H'''(\alpha)>0$ for $\alpha\in[\frac{1}{3}, 1)$, which implies $H''(\alpha)$ is increasing. Also $H''(1)<0$. Hence $H''(\alpha)<0$ in $\alpha\in[\frac{1}{3}, 1)$ and since $H'(1)>0$, which means that $H'(\alpha)>0$ for any $\alpha\in[\frac{1}{3},1)$. Since $H(\frac{1}{3})>0$ so we conclude that $H(\alpha)>0$ in $[\frac{1}{3}, 1)$. Therefore $M_{1}(\alpha)\geq0$ in $\alpha\in[\frac{1}{3},1)$, which further gives $M(q,\alpha)\geq0$ for all even $k\geq8$ and $\alpha\in[\frac{1}{3},1)$. On the other hand, $I_{3}>0$ and hence  $C_{k}(q,\alpha)>0$ for all even $k\geq8$ and $\alpha\in[\frac{1}{3},1)$. Now for the two remaining even numbers $k=4$, and for $k=6$, $C_{k}(q,\alpha)$ has the following expressions
\begin{align*}
C_{4}(q,\alpha)&=\frac{\gamma_{1}}{4!}\frac{q_{n}^{4}}{Q_{n}^{4}}\displaystyle\prod_{i=1}^{3}(i-\gamma_{1})+
\frac{\gamma_{2}}{4!}\frac{q_{n}q_{n+1}^{3}}{Q_{n}^{4}}\displaystyle\prod_{i=1}^{3}
(i-\gamma_{2})-\frac{\alpha}{4!}\frac{q_{n}q_{n+1}^{3}}{Q_{n}^{4}}\displaystyle\prod_{i=1}^{3}(i-\alpha),\\
C_{6}(q,\alpha)&=\frac{\gamma_{1}}{6!}\frac{q_{n}^{6}}{Q_{n}^{6}}\displaystyle\prod_{i=1}^{5}(i-\gamma_{1})+
\frac{\gamma_{2}}{6!}\frac{q_{n}q_{n+1}^{5}}{Q_{n}^{6}}\displaystyle\prod_{i=1}^{5}
(i-\gamma_{2})-\frac{\alpha}{6!}\frac{q_{n}q_{n+1}^{5}}{Q_{n}^{6}}\displaystyle\prod_{i=1}^{5}(i-\alpha).
\end{align*}
Since $q_{n}\geq q_{n+1}$ for $n\in\mathbb{N}$ so one can write
\begin{align}{\label{c-4}}
C_{4}(q,\alpha)&\geq J_{1}(\alpha)\frac{q_{n}q_{n+1}^{3}}{Q_{n}^{4}},
\end{align} where we denote
\begin{align*}
J_{1}(\alpha)=&\frac{\gamma_{1}}{4!}\displaystyle\prod_{i=1}^{3}(i-\gamma_{1})+\frac{\gamma_{2}}{4!}\displaystyle
\prod_{i=1}^{3}(i-\gamma_{2})-\frac{\alpha}{4!}\displaystyle\prod_{i=1}^{3}(i-\alpha).
\end{align*}Plugging the values of $\gamma_{1}=\frac{1-\alpha}{2}$, $\gamma_{2}=\frac{1+\alpha}{2}$ in $J_{1}(\alpha)$ we get a simplified form as below
\begin{align*}
J_{1}(\alpha)=\frac{1}{4!8}(\alpha^2-6\alpha+5)(7\alpha^2-6\alpha+3).
\end{align*} It is clear that $J_1(\alpha)>0$ in $[\frac{1}{3}, 1)$ which shows that $C_{4}(q,\alpha)>0$ in $[\frac{1}{3}, 1)$.\\
Similarly one has
\begin{align}
C_{6}(q,\alpha)&\geq J_{2}(\alpha)\frac{q_{n}q_{n+1}^{5}}{Q_{n}^{6}},
\end{align}
where we denote
\begin{align*}
J_{2}(\alpha)=&\frac{\gamma_{1}}{6!}\displaystyle\prod_{i=1}^{5}(i-\gamma_{1})+\frac{\gamma_{2}}{6!}\displaystyle\prod_{i=1}^{5}
(i-\gamma_{2})-\frac{\alpha}{6!}\displaystyle\prod_{i=1}^{5}(i-\alpha).
\end{align*}Inserting the values of $\gamma_{1}$ and $\gamma_{2}$ in $J_{2}(\alpha)$, we obtain
\begin{align*}
J_{2}(\alpha)=&\frac{1}{6!32}(31\alpha^{6}-480\alpha^{5}+2515\alpha^{4}-7200\alpha^{3}+8029\alpha^{2}-3840\alpha+945).
\end{align*}Differentiating $J_{2}(\alpha)$ up to four times, we get
\begin{align*}
J'_{2}(\alpha)&=\frac{1}{6!32}(186\alpha^{5}-2400\alpha^{4}+10060\alpha^{3}-21600\alpha^{2}+16058\alpha-3840),\\
J''_{2}(\alpha)&=\frac{1}{6!32}(930\alpha^{4}-9600\alpha^{3}+30180\alpha^{2}-43200\alpha+16058),\\
J'''_{2}(\alpha)&=\frac{1}{6!32}(3720\alpha^{3}-28800\alpha^{2}+60360\alpha-43200),\\
J''''_{2}(\alpha)&=\frac{1}{192}(93\alpha^{2}-480\alpha+503).
\end{align*}\\
Note that $J''''_{2}(\alpha)>0$ in $[\frac{1}{3}, 1)$ which implies that $J'''_{2}(\alpha)$ is increasing in $[\frac{1}{3}, 1)$ but $J'''_{2}(1)<0$, which says that $J''_{2}(\alpha)$ is decreasing in $[\frac{1}{3}, 1)$ but $J''_{2}(\frac{27}{50})>0$ implies that $J''_{2}(\alpha)>0$ in $[\frac{1}{3}, \frac{27}{50}]$. That is $J'_{2}(\alpha)$ is increasing in $[\frac{1}{3}, \frac{27}{50}]$ but $J'_{2}(\frac{1}{3})>0$ and $J_{2}(\frac{1}{3})>0$, which finally gives $J_{2}(\alpha)>0$ in $[\frac{1}{3}, \frac{27}{50}]$. One can also write $J_{2}(\alpha)$ as $J_{2}(\alpha)=\frac{(9-\alpha)f(\alpha)}{6!32}$,
where $f(\alpha)=-31\alpha^{5}+201\alpha^{4}-706\alpha^{3}+846\alpha^{2}-415\alpha+105$. Then
\begin{align*}
 f'(\alpha)&=-155\alpha^{4}+804\alpha^{3}-2118\alpha^{2}+1692\alpha-415<-6(25\alpha^{4}-134\alpha^{3}+353\alpha^{2}-282\alpha+69).
 \end{align*}We observe that $25\alpha^{4}-134\alpha^{3}+353\alpha^{2}-282\alpha+69>0$ in $\alpha\in[27/50, 1)$. Therefore $f'(\alpha)<0$. Also since $f(1)=0$, so
 $f(\alpha)>0$ in $\alpha\in[27/50, 1)$. Consequently $J_{2}(\alpha)>0$ $\forall\alpha\in[27/50,1)$, and hence $J_{2}(\alpha)>0$ for $\alpha\in[\frac{1}{3}, 1)$. This gives $C_{6}(q,\alpha)>0$ for $\alpha\in[\frac{1}{3}, 1)$. Hence we prove the Lemma.
\end{proof}

\end{document}
